% LSPV-4 REPORT::MODELBRIDGE-0027::2026-10-12T14:58:27Z
% !TEX program = xelatex
\documentclass[11pt]{article}
\usepackage[margin=0.88in]{geometry}
\usepackage{fontspec,amsmath,amssymb,amsthm,booktabs,longtable,array,enumitem,xurl,hyperref}
\setlength{\emergencystretch}{3em}
\newtheorem{theorem}{Theorem}
\newtheorem{corollary}[theorem]{Corollary}
\newcommand{\EXACT}{\textbf{PROVED/EXACT}}
\newcommand{\KNOWN}{\textbf{KNOWN THEOREM}}
\newcommand{\COMP}{\textbf{COMPUTATIONAL}}
\newcommand{\CAND}{\textbf{CANDIDATE}}
\newcommand{\OPEN}{\textbf{OPEN}}
\newcommand{\RETIRED}{\textbf{FALSIFIED/RETIRED}}
\newcommand{\same}{\textsc{same\_realization}}
\title{MODELBRIDGE-0027:\\Bare Full-Trace Factorization Leaks the Source and Supplies No SAT Lower Bound}
\author{LSPV-4 --- Universal Core Bridge}
\date{12 October 2026}

\begin{document}
\sloppy
\maketitle

\begin{abstract}
The newest MAIN synthesis reduces program-specific checkpoint schedules to
full-trace observation and asks whether clean-core compatibility factors
through the same realized trace.  This report proves a sharp qualification.
In the standard machine model, the full trace contains the initial
configuration and therefore the encoded formula.  Any pointwise compatibility
service computable from the formula under a hypothetical polynomial SAT
decider consequently factors through the trace by an extractor that ignores
all later solver states and performs separately charged computation.  Its cost
retains the arbitrary solver exponent.  If the target is instead the explicit
full compatibility table, output length is exponential but that is only a
representation/output lower bound, not a lower bound for one-bit SAT decision.
Thus \texttt{BARE-FULL-TRACE-FACTORIZATION-AS-BRIDGE1} is retired.  A
nontrivial internally derived trace projection or a direct lower bound for all
algorithms deciding the fixed NP-complete clean core remains open.  P versus
NP remains open.
\end{abstract}

\section{Executive result and claim card}
\begin{center}\footnotesize
\begin{tabular}{p{0.43\linewidth}p{0.26\linewidth}p{0.23\linewidth}}
\toprule
Item & Status & Reason\\
\midrule
\texttt{INPUT-TRACE-SOURCE-LEAKAGE1} & \EXACT & Initial trace state decodes the source\\
Finite hostile-control audit & \COMP & 104,805 assertions, zero failures\\
\texttt{BARE-FULL-TRACE-FACTORIZATION-AS-BRIDGE1} & \RETIRED & Extractor may ignore the computation\\
\texttt{INTERNAL-DERIVED-TRACE-OR-DIRECT-CORE-LB1} & \OPEN & No universal transfer or lower bound\\
P versus NP & \OPEN & Load-bearing arrow remains\\
\bottomrule
\end{tabular}
\end{center}

\paragraph{Claim card.}
The status before this run was \OPEN.  The attacked implication was
\[
 \forall A\in\mathsf{PolyDec}(K)\ \exists E_A\ \forall F,\alpha:
 E_A(\operatorname{Tr}_A(F),\alpha)=H(F,\alpha)
 \quad\Longrightarrow\quad
 A(F)\text{ exposed or paid for }H.                       \tag{1}
\]
Here $K$ is the canonical duplicate-free proper signed Monotone
$3$SAT-$(2,2)$ clean-core language, $N=|\operatorname{enc}(F)|$, $A$ is a
uniform deterministic exact decider, and $\operatorname{Tr}_A(F)$ is the
complete exact-clock configuration trace including time zero.  The resource
measure is joint time/space of one declared ordinary realization.  The
\same{} object is the run of $A$ on $F$.  A single counterexample extractor
that ignores every state after time zero falsifies the claimed implication.

\section{Exact compatibility variants}
For a consistent partial assignment $\alpha\in\{0,1,\star\}^n$, distinguish
the pointwise service
\[
 h(F,\alpha)=1
 \quad\Longleftrightarrow\quad
 \exists\beta\supseteq\alpha\;[\beta\models F]             \tag{2}
\]
from the explicit compatibility relation
\[
 \mathcal H_F=\{(\alpha,h(F,\alpha)):
                 \alpha\in\{0,1,\star\}^n\}.              \tag{3}
\]
The first returns one bit.  The second has $3^n$ indexed entries before any
compression.  Conflating them makes the transfer target ambiguous.

The clean-core promise is not preserved by restriction.  Therefore a query
to a clean-core decider must use a polynomial reduction
$r(F\upharpoonright\alpha)$ back to $K$, as established in
MODELBRIDGE-0025.  This adapter is part of the extractor cost.

\section{INPUT-TRACE-SOURCE-LEAKAGE1}
\begin{theorem}[Source leakage from the full trace]
Let $M$ be any standard deterministic machine whose initial configuration
$C_0(x)$ contains a decodable self-delimiting input string $x$.  Let
\[
 \operatorname{Tr}_M(x)=(C_0(x),C_1(x),\ldots,C_{T_M(x)}(x)).
\]
For every partial function $g$ and every algorithm $B$ computing $g(x)$,
there is an extractor $E_{M,B}$ such that
\[
 E_{M,B}(\operatorname{Tr}_M(x))=g(x),                     \tag{4}
\]
and $E_{M,B}$ reads no configuration after $C_0(x)$ except to parse the trace
container.  Its work is the input-decoding cost plus the work of $B$; that
work is not work performed by the original execution $M(x)$.
\end{theorem}

\begin{proof}
Read $C_0(x)$, apply the fixed decoder for the machine/input convention to
recover $x$, and run $B(x)$.  Equation (4) follows by correctness of the
decoder and $B$.  The construction never inspects the transition from $C_0$
to $C_1$, so it cannot certify that the original computation derived,
materialized, retained, or paid for $g(x)$.  The extractor is a second
execution with its own time and space. \qed
\end{proof}

\begin{corollary}[Pointwise compatibility factorization is vacuous under
$K\in P$]
Assume the clean-core language $K$ has a uniform deterministic decider $A$
with time $m^{k+o(1)}$ and space $m^{d+o(1)}$.  If the restriction adapter has
length $|r(F\upharpoonright\alpha)|\le N^{\lambda+o(1)}$, construction time
$N^{u+o(1)}$, and workspace $N^{v+o(1)}$, then the extractor that uses only
the initial trace state and calls $A$ satisfies
\[
 T_E(N)\le N^{\max\{u,\lambda k\}+o(1)},\qquad
 S_E(N)\le N^{\max\{v,\lambda d\}+o(1)}.                  \tag{5}
\]
With $N^b$ adaptive queries, the time exponent includes $b+\lambda k$.
Thus pointwise $h(F,\alpha)$ factors through the full trace whenever $K\in P$,
but the arbitrary exponents of the assumed decider remain and no fixed
forbidden endpoint is reached.
\end{corollary}

\begin{proof}
Decode $F$ from $C_0(F)$, construct $r(F\upharpoonright\alpha)$, and invoke
$A$.  Charge the adapter and the invoked run.  Taking the maximum of the
polynomial exponents gives (5); query multiplicity adds $b$. \qed
\end{proof}

\begin{corollary}[Explicit-table output does not lower-bound decision]
Any literal serialization of (3) writes at least $3^n$ answer positions.
This is an output-size lower bound for that representation.  It does not imply
that deciding the single bit $h(F,\star^n)$ takes superpolynomial time or
space, because an arbitrary decider need not enumerate (3).
\end{corollary}

\section{Falsification and weakest repair}
The source-leakage theorem falsifies (1).  Bare existence of
$E_A(\operatorname{Tr}_A(F),\alpha)$ is too weak: it can hold solely because
the trace transports the original input.  The factorization statement must
therefore specify at least one of the following nontrivial repairs.

\begin{enumerate}[leftmargin=*]
\item \textbf{Derived-state repair:} the extractor receives only a declared
projection of states after input erasure and cannot recover the formula by a
fixed decoder.  This immediately reopens wrapper and alternate-representation
attacks; no implementation-invariant projection is known.
\item \textbf{Internal-cost repair:} the resource inequality is proved for the
same deciding execution, without running $A$ again or invoking an external
restriction adapter.  This is the genuine \same{} target and remains \OPEN.
\item \textbf{Direct-core repair:} bypass semantic exposure and prove that no
uniform deterministic polynomial-time algorithm decides the fixed NP-complete
clean-core language.  This would suffice for $P\ne NP$ without transfer to
unrestricted SAT, but it is exactly the open lower-bound objective.
\end{enumerate}

The weakest useful next theorem is therefore not bare factorization but a
cost-preserving internal trace inequality that survives decoding, wrappers,
and program-specific state encodings, or a direct all-algorithm lower bound on
$K$.

\section{Bounded reproducible audit}
The script
\path{experiments/review/full_trace_source_leakage_audit.py} exhaustively
tested 9 finite families, 5 encodings, and 6 wrappers, for 270 traces.  It
decoded the source solely from time zero, compared pointwise compatibility on
all ternary partial assignments, and verified that later trace suffixes differ
across DIRECT, DEAD, COPY, DELAY, INTERP, and RECOMP presentations.  All
104,805 assertions passed; 104,490 point queries and 104,490 explicit table
entries were checked in 2.121457 seconds.  The raw result is
\path{results/raw/review/2026-10-12_full_trace_source_leakage.json}.

This audit is \COMP.  It checks finite identities only; the symbolic theorem
above is what establishes the general source-leakage result.  The invocation
was one bounded chunk and completed below the repository's 30-second limit.

\section{Hostile controls and alternate representations}
\begin{longtable}{p{0.30\linewidth}p{0.61\linewidth}}
\toprule
Control & Outcome\\
\midrule
Horn, 2SAT, affine/Tseitin, bounded treewidth & Source decoding and external point queries work even where compatibility is easy; no hardness signal is isolated.\\
Independent products & Point queries stream; no answer vector or additive peak space is forced.\\
Padding and redundant clauses & Initial input size and adapter work are charged; semantic redundancy does not become solver-paid exposure.\\
DIRECT, DEAD, COPY, DELAY & Later traces and resource use change while the time-zero source remains decodable.\\
INTERP/RECOMP & Alternate internal presentations do not affect source extraction from $C_0$.\\
Circuit-envelope encoding & A fixed decoder recovers the source; encoding changes do not make later states necessary.\\
One-bit decision output & Only $h(F,\star^n)$ is required; explicit $3^n$ output is optional work.\\
\bottomrule
\end{longtable}

\section{Primary-literature preflight and model matching}
The major obstacles are (i) separating information merely transported in the
input from information internally derived by the computation, and (ii)
connecting any derived trace property to an unconditional lower bound in the
same machine model.

\begin{itemize}[leftmargin=*]
\item Cook's SAT completeness theorem is a \KNOWN{} direct transfer for
language membership.  It justifies the reduction adapter and the fact that a
direct non-$P$ proof for one NP-complete core would finish the objective.  It
does not provide the missing lower bound.
\item Pippenger--Fischer oblivious simulation is a \KNOWN{} theorem about
simulating access patterns with $O(T\log T)$ time.  It does not distinguish
input-carried from derived semantics or force $H$ into mutable state.
\item Fortnow--Lipton--van Melkebeek--Viglas supply \KNOWN{} SAT time--space
lower bounds in stated random-access/small-space regimes.  Equation (5) retains
the arbitrary exponent $k$, so no model-matched landing follows.
\item Baker--Gill--Solovay and Aaronson--Wigderson are \KNOWN{} barrier results.
Source decoding and external SAT calls relativize and algebrize; this theorem
is a guard, not a barrier escape.
\end{itemize}

\paragraph{Cross-domain transfer.}
Control-theoretic observability assumes a fixed state/output map; here the
time-zero output already contains the input and equivalent programs change
later state coordinates.  It is a \textbf{FAILED TRANSFER} as an unrestricted
lower-bound theorem.  Information-flow taint or provenance systems exactly
distinguish source carriage from derivation only relative to a declared
operational semantics; applying them to arbitrary equivalent Turing/RAM
programs is an \textbf{ANALOGY/CANDIDATE} needing an invariant compiler.
VLSI locality, proof-system interpolation, and communication cuts remain
restricted-model analogies without an all-algorithm simulation theorem.

\section{Relativization, algebrization, Natural Proofs, and uniformity}
\paragraph{Relativization.}
Give the extractor and decider the same oracle.  Recovering $F$ from $C_0$ and
issuing oracle-relative restriction queries still works.  No line separates
$P^O$ from $NP^O$.

\paragraph{Algebrization.}
Replacing an oracle by its low-degree extension does not attribute the
extractor's second computation to the original run.  The result survives and
supplies no nonalgebrizing hinge.

\paragraph{Natural Proofs.}
No constructive large circuit property is asserted.  If a future trace
projection is promoted to a circuit property, constructivity, largeness,
usefulness, and the exact circuit class must be audited.

\paragraph{Uniformity and quantifiers.}
The extractor is allowed to depend on the fixed program $A$, but one finite
$E_A$ must work for all input lengths.  Equation (5) uses
$\forall A\,\exists k,d,E_A$; it does not assume a common polynomial exponent.
Hardwiring an answer table per length is not allowed.

\section{Post-quantum technique ledger}
\begin{itemize}[leftmargin=*]
\item Schur squares, quadratic relation spaces, Hilbert functions, syzygies,
and graded Betti numbers: \textbf{FAILED TRANSFER} for bare full-trace
factorization.  A source decoder can reconstruct and analyze any chosen
algebra externally; no such algebra is forced into solver state.
\item Collision entropy: \textbf{EXACT OBJECT/THEOREM TRANSFER} only for a
fixed map $L$, with
$|\operatorname{im}L|\ge |A|^2/\operatorname{Coll}(L)$.  It lower-bounds an
image, not the original decision cost.
\item LWE-style black-box extraction and special-soundness collision
extraction: \textbf{ANALOGY/CANDIDATE}.  They require a new cost-preserving
compiler; cryptographic assumptions cannot be the unconditional endpoint.
\item Hash/random-oracle binding: \textbf{FAILED TRANSFER} unless all hardness
and random-oracle assumptions are removed.
\end{itemize}

\section{CORPUS REVIEWED}
This run surveyed the current GitHub corpus and read root \path{main.tex},
\path{AGENTS.md}, \path{CODEX.md}, the pnp-research skill and protocol
references, all current task specifications, the newest verified MAIN
SYNTHESIS-0024, the actual \path{main4.tex}/\path{main5.tex}, and the newest
verified LSPV-1 through LSPV-4 reports.  Materially relevant historical reports
included HEVALLB1, APIFLAT1, MODELBRIDGE-0022 through MODELBRIDGE-0026, CUTGAP1,
PRODSTREAM1, and EXOGATE1.  Repository searches covered full trace, input
configuration, compatibility, restriction adapter, source factorization,
wrappers, and the HEC/VSC-HEC/G-PPC-HEC/CRS-PPC/CCSPC causal branch.  Relevant
governance notes and cited experiment/result artifacts were inspected.

The predecessor LSPV-3 artifact
\path{research/agents/lspv-3/2026-10-12_2338_LSPV-3_EXOGATE1.tex} is
substantive and has the required header and final sentinel.  The requested
recovery path for MODELBRIDGE-0022 now exists canonically at
\path{research/agents/lspv-4/2026-10-12_1002_LSPV-4_MODELBRIDGE-0022.tex}; it
was read back successfully, so no duplicate recovery file was created.

\section{DUPLICATION / CONTRADICTION AUDIT}
MODELBRIDGE-0022 already proves that black-box extraction is a separately
charged computation.  MODELBRIDGE-0023 and -0025 already retain the arbitrary
solver exponent and charge the clean-core restriction adapter.  The new
increment is to apply those guards to SYNTHESIS-0024's exact full-trace target:
because a full trace contains $C_0(F)$, even a nominally ``trace-only''
extractor can be source-only.  This narrows, rather than contradicts, the
synthesis.  No earlier verified theorem established internal full-trace
factorization.

\section{WHAT CHANGED RELATIVE TO ROOT MAIN}
Root MAIN correctly marks full-trace semantic factorization \OPEN.  This
report shows that the bare equation
$E_A(\operatorname{Tr}_A(F))=h(F,\alpha)$ is not a sufficient bridge, because
it is satisfied by time-zero source decoding plus external computation.  The
living LSPV-4 region should distinguish source-carried factorization from an
internally derived, cost-preserving \same{} trace inequality.

The actual \path{main4.tex} remains an achievable-region framework whose
SAT-specific contraction is open.  The actual \path{main5.tex} remains a
conditional WRAPPER composition manuscript whose all-input cheap acquisition
contract is open.  Neither is redefined as a full-trace lower bound.

\section{ADOPTED / REJECTED / RETIRED PRIOR COMPONENTS}
\begin{itemize}[leftmargin=*]
\item \textbf{ADOPTED:} exact-clock full traces; corrected clean core;
restriction adapters; wrapper controls; explicit all-cost accounting; direct
clean-core non-$P$ as a sufficient final target.
\item \textbf{REJECTED:} treating input carriage as semantic derivation;
uncharged calls by $E_A$; transferring explicit compatibility-table output
size to one-bit decision time.
\item \textbf{FALSIFIED/RETIRED:}
\texttt{BARE-FULL-TRACE-FACTORIZATION-AS-BRIDGE1}.
\item \textbf{OPEN:} an implementation-invariant derived trace projection; an
internal \same{} cost inequality; a direct all-algorithm lower bound for $K$;
a model-matched unconditional contradiction.
\end{itemize}

\section{STALE OR OVERSTATED CLAIMS FOUND}
Any wording that treats ``computable from the full trace'' as synonymous with
``computed, exposed, or paid for by the deciding execution'' is overstated.
Any wording that treats the $3^n$ explicit compatibility table as necessary
output of SAT decision is likewise overstated.  These must be read as
\RETIRED{} unless a derived-state and same-realization theorem is supplied.

\section{Precise next target}
Define a projection $D_A(F)$ of the realized trace that excludes the decodable
input carrier yet is invariant under polynomial-delay, copying,
recomputation, interpretation, and alternate state encodings.  Then prove one
of the following with all costs charged:
\[
 \forall A\in\mathsf{PolyDec}(K)\ \exists E_A\ \forall F,\alpha:
 E_A(D_A(F),\alpha)=h(F,\alpha),                          \tag{6}
\]
where (6) yields a resource inequality for the same execution, or a direct
superpolynomial lower bound for every uniform deterministic decider of $K$.
If no nontrivial invariant $D_A$ survives the wrapper tests, retire the trace
factorization branch and concentrate on the direct clean-core/constructive
WRAPPER routes.

\begin{thebibliography}{9}
\bibitem{Cook71} S. Cook, ``The complexity of theorem-proving procedures,''
STOC 1971, 151--158, \url{https://doi.org/10.1145/800157.805047}.
\bibitem{PF79} N. Pippenger and M. J. Fischer, ``Relations among complexity
measures,'' \emph{J. ACM} 26(2), 361--381, 1979,
\url{https://doi.org/10.1145/322123.322138}.
\bibitem{BGS75} T. Baker, J. Gill, and R. Solovay, ``Relativizations of the
$P=?NP$ question,'' \emph{SIAM J. Comput.} 4(4), 431--442, 1975,
\url{https://doi.org/10.1137/0204037}.
\bibitem{AW09} S. Aaronson and A. Wigderson, ``Algebrization: A New Barrier in
Complexity Theory,'' \emph{ACM TOCT} 1(1), 2009, ECCC TR08-005.
\bibitem{FLMV05} L. Fortnow, R. Lipton, D. van Melkebeek, and A. Viglas,
``Time--space lower bounds for satisfiability,'' \emph{J. ACM} 52(6),
835--865, 2005, \url{https://doi.org/10.1145/1101821.1101822}.
\bibitem{DD21} A. Darmann and J. D\"ocker, ``On simplified NP-complete
variants of Monotone 3-Sat,'' \emph{Discrete Applied Mathematics} 292,
45--58, 2021, \url{https://doi.org/10.1016/j.dam.2020.12.010}.
\end{thebibliography}

\noindent\textbf{Final verification sentinel:} report body complete; strict
statuses preserved; P versus NP remains \OPEN.
% END_OF_VERIFIED_REPORT::LSPV-4::MODELBRIDGE-0027
\end{document}
